\documentclass[11pt,a4paper]{article}
\usepackage[utf8]{inputenc}
\usepackage{amsmath,amssymb,amsfonts}
\usepackage{geometry}
\usepackage{hyperref}
\geometry{margin=1in}

\title{\textbf{Auto Release Control (ARC) Dynamic Peak Limiting \& Gain Smoothing in Time-Dilated Audio Engines}}
\author{Time Dilation DAW Research Group --- Version 0.0.1}
\date{August 2026}

\begin{document}
\maketitle

\section{Overview}
In relativistic time-dilated audio synthesis, dynamic time warping and feedback cascades can produce transient amplitude spikes exceeding standard $0.0 \text{ dBFS}$ audio ceilings. To prevent output clipping and digital inter-sample distortion while maintaining acoustic warmth, the \texttt{out\textasciitilde} (\texttt{OutNode}) master object incorporates an automatic \textbf{ARC Safety Peak Limiter} operating at a ceiling of $-1.5 \text{ dBFS}$.

\section{Ceiling \& Gain Threshold Math}
The linear ceiling threshold $C$ for a peak ceiling of $-1.5 \text{ dBFS}$ is defined by:
\begin{equation}
C = 10^{\frac{-1.5}{20}} \approx 0.8413951
\end{equation}

For each audio sample $s \in [0, N-1]$ in a stereo block, the instantaneous peak magnitude $p[s]$ is computed as:
\begin{equation}
p[s] = \max\left(|s_L[s]|, |s_R[s]|\right)
\end{equation}

The target gain multiplier $g_{\text{target}}[s]$ required to constrain $p[s]$ within the ceiling $C$ is:
\begin{equation}
g_{\text{target}}[s] = \begin{cases}
1.0 & \text{if } p[s] \le C \\
\frac{C}{p[s]} & \text{if } p[s] > C
\end{cases}
\end{equation}

\section{Two-Stage Attack \& Adaptive ARC Release Filter}
To suppress transient clipping instantly without introduced harmonic distortion or audible pumping during sustained signals, the limiter uses an \textbf{Auto Release Control (ARC)} mechanism.

\subsection{Attack Phase}
The attack filter operates at an ultra-fast time constant $\tau_{\text{att}} = 0.05 \text{ ms}$ ($50 \ \mu\text{s}$):
\begin{equation}
\alpha_{\text{att}} = 1.0 - \exp\left(-\frac{1.0}{f_s \cdot 0.00005}\right)
\end{equation}

When $g_{\text{target}}[s] < g[s-1]$, the limiter triggers attack attenuation:
\begin{equation}
g[s] = g[s-1] + \alpha_{\text{att}} \cdot (g_{\text{target}}[s] - g[s-1])
\end{equation}

\subsection{Adaptive ARC Release Phase}
When $g_{\text{target}}[s] \ge g[s-1]$, the release time constant dynamically adjusts based on the transient intensity ratio $R = p[s] / C$:
\begin{equation}
\tau_{\text{rel}}(R) = \begin{cases}
10 \text{ ms} & \text{if } R < 1.5 \quad (\text{Fast recovery for transient spikes}) \\
300 \text{ ms} & \text{if } R \ge 1.5 \quad (\text{Smooth recovery for sustained peaks})
\end{cases}
\end{equation}

The release smoothing coefficient $\alpha_{\text{rel}}$ is computed sample-by-sample:
\begin{equation}
\alpha_{\text{rel}} = 1.0 - \exp\left(-\frac{1.0}{f_s \cdot (\tau_{\text{rel}} / 1000.0)}\right)
\end{equation}
\begin{equation}
g[s] = g[s-1] + \alpha_{\text{rel}} \cdot (g_{\text{target}}[s] - g[s-1])
\end{equation}

\section{Output Signal Scaling}
The final limited stereo sample pair $(y_L[s], y_R[s])$ delivered to the DAW output buffer is:
\begin{equation}
y_L[s] = s_L[s] \cdot g[s], \quad y_R[s] = s_R[s] \cdot g[s]
\end{equation}
This guarantees zero hard clipping at $\pm 1.0$, complete protection of downstream audio devices, and transparent ARC release dynamics.

\end{document}
